\labeledsubsection{Full experiment}
With the previous sub-section the plugin-system, with the \texttt{agent}/\lstinline{instrumentation}, is defined and a plugin exists, is loaded and the task-method was invoked.

The full experiment goes a bit more into detail and uses cleaner code.
Overall, the experiment not only has a task-method, but also a enable- and disable-method.
All plugin information (such as the plugin-class, \lstinline{@PluginInfo} and all found methods) are stored in a separate data-class.
Said data-class also implements helper-functions to make invoking said methods easier by simply calling \lstinline{Plugin::enable}, \lstinline{Plugin::disable} or \lstinline{Plugin::task}.
Additionally, certain constraints, such as if the plugin file exists, are implemented and multiple plugins can be loaded by defining and searching in a specific folder for JAR files.

The full source code can be found in the appendix starting from page~\pageref{subsec:Experiment Source Code}.

\input{tex/003_experiment/001_fulll_experiment/001_potential_risks_plugin_system.tex}
\clearpage
\labeledsubsubsection{Potential risks and threads extended scope}
The previous sub-sub-section focused on a \textit{plugin-system} /\\ a \underline{planned dependency-injection}.
However, the entire scope can be extended to also include \underline{unwanted/unplanned dependency-injections} and other similar attacks.

The main problem with most applications is that they use external dependencies.
Such libraries can be hijacked and modified, to then either include malicious code or return wrong results.
The linux operating system solves this issue by storing all binaries and shared libraries at a certain location and \underline{only} the \texttt{root} user has (/ should have) access to it.
This way, no user, nor application, can inject, hijack or modify anything.
Other operating systems, like Windows, do not have this level of security.
Only the system libraries are securely stored, most other libraries and binaries are usually accessible by the current user and any application that is run by them.
By utilizing and properly setting up tools like UAC\footnote{User-Access-Control} it is possible to archive a similar security level.

Another huge risk are \textit{mods} for games.
Many function on a similar base as the \textit{plugin-system}.
Thus, any mod could, potentially, be malicious.
This is especially the case when a user has to mod the game themselves, instead of letting a tool (such as a game-launcher) do this.
By giving the option for practically everyone to create a mod and upload it anywhere, for others to download and install it, the risk of someone publishing a malicious mod is high.
There is also no way of trusting user-statistics, comments and other indications on third-party sides.
\labeledsubsubsection{How to improve security?}
To overall increase security and protect against dependency-injection in any form, one must first of all secure their own systems.
Only proper permissions, file access and access-control will lead to a safe system.

Next, it is important to inform users of \textit{plugin-systems} or \textit{mods} that not every \textit{plugin}/\textit{mod} has necessarily good intensions.
Creating a central place for all \textit{mods} or \textit{plugins} to be uploaded to also increases the security.
Such a place must require a developer to sign-up so that, in case a malicious artifact was found, this user can be banned, the artifact removed and users informed about it.
Thus, a platform also has to have a comment or ranking system, including the option to report malicious uploads.\\
Even better, if the application has an integrated way of installing a \textit{plugin} or \textit{mod} in either the application itself or a launcher.
A good example for this is \texttt{Visual Studio Code}, which features an integrated and controlled marketplace for all sorts of \textit{plugins}.
Another example would be \texttt{NetBeans}, \texttt{IntelliJ IDEA/Jetbrains-Applications} or \texttt{Google Chrome/Chromium}.
All feature such a central hub for uploading and installing \textit{plugins}.

Finally, to be truly safe from dependency-injections one can execute an application only in an isolated, maybe containerized, environment with minimal to non additional tools and permissions.

Additionally, it may be possible to scan \textit{plugins} or \textit{mods} for malicious or unwanted code.
However, this process would be quiet expensive and probably also only works if the source code is available.